\documentclass[doublespace,prb,aps]{article}
\usepackage[utf8]{inputenc}
\usepackage{amsmath}
\usepackage{lipsum} % to generate some filler text
\usepackage{fullpage}
\usepackage{mathrsfs}
\usepackage{svg}
\usepackage{xr}
\usepackage{amsthm,url,graphicx,booktabs,lipsum,color,bm,caption,subcaption,soul}
\usepackage{pifont,tikz,paralist,multirow,amssymb}
\usepackage{chemarr,natbib}
\usepackage{xifthen}
\usepackage{enumerate}
\usepackage{titlesec}
\newcounter{reviewer}
\setcounter{reviewer}{0}
\newcounter{point}[reviewer]
\setcounter{point}{0}

% This refines the format of how the reviewer/point reference will appear.
%\renewcommand{\thepoint}{C\,\thereviewer.\arabic{point}}

\renewcommand{\thepoint}{Comment\,\thereviewer.\arabic{point}}

% command declarations for reviewer points and our responses

\newenvironment{point1}
   {\refstepcounter{point} \bigskip \noindent {\textbf{Associate Editor Comment~}} ---\ }
   {\par }

\newcommand{\reviewersection}{\stepcounter{reviewer} \bigskip \hrule
                  \section*{Reviewer \thereviewer}}

% \newenvironment{point}

%    {\refstepcounter{point} \bigskip \noindent {\textbf{Reviewer~\thepoint} } ---\ }
%    {\par }

\newenvironment{point}
   {\refstepcounter{point} \bigskip \noindent {\textbf{Reviewer~\thepoint} } ---\ }
   {\par }


\newcommand{\shortpoint}[1]{\refstepcounter{point}   \bigskip \noindent 
	{\textbf{Reviewer~Point~\thepoint} } ---~#1  }

\newenvironment{reply}
   {\medskip \noindent \color{black} \begin{sf}\textbf{Reply}:\  }
   {\medskip \par \end{sf}}

\newenvironment{revision}
   {\medskip \noindent \color{blue} \begin{sf}\textbf{Revision}:\textit{\ }}
   {\medskip \par \end{sf}}


\newcommand{\shortreply}[2][]{\medskip \noindent \begin{sf}\textbf{Reply}:\  #2
	\ifthenelse{\equal{#1}{}}{}{ \hfill \footnotesize (#1)}%
	\medskip \end{sf}}

\makeatletter
\newcommand*{\addFileDependency}[1]{% argument=file name and extension
  \typeout{(#1)}
  \@addtofilelist{#1}
  \IfFileExists{#1}{}{\typeout{No file #1.}}
}
\makeatother

\newcommand*{\myexternaldocument}[1]{%
    \externaldocument{#1}%
    \addFileDependency{#1.tex}%
    \addFileDependency{#1.aux}%
}

\myexternaldocument{main}
%\externaldocument[I-]{main}


\def\thesubsubsection{\arabic{subsubsection}}
\titleformat{\subsubsection}[runin]
{\normalfont\bfseries}
{\indent\thesubsubsection)}{0.5em}{}


\begin{document}
\section*{Notes on revision made to manuscript Paper $\#$ct-2026-001787: \\ Improving binding free energy predictions with Swap Monte Carlo for water sampling 
}
\vspace{10pt}

The authors sincerely thank the editor and all reviewers for their careful reading of the manuscript and for providing thoughtful, constructive comments. 
We have carefully considered each point and revised the manuscript accordingly. 
Significant changes in the revised manuscript are highlighted in \textcolor{blue}{blue} for the reviewers' convenience. 
Our point-by-point responses are provided below.
% Let's start point-by-point with Reviewer 1

%\section*{Response to the Associate Editor}
% General intro text goes here
%\begin{point1}
%For your work, I have received three independent review reports from experts 
%\end{point1}


%\begin{reply}
%Thank you so much for your constructive comments. 

%\end{reply}

\section*{Response to the reviewers}


%%%%%%%%%%%%%%%%%%%%%%%%%%%%%%%%%%%%%%%%%%%%%%%%%%%
\reviewersection
This is an interesting MC study on water sampling in the context of MD simulations/free energy calculations, 
and I believe it will be of significant interest to workers in the area. 
However, there are several aspects I would recommend improving.

% Point one description 
\begin{reply}
We thank the reviewer for the positive assessment and for the detailed suggestions, 
which have helped strengthen the manuscript.
\end{reply}

%%%%%%%%%%%%%%%%
\begin{point}
Some critical details are missing in methods; in general, the authors should provide enough detail someone should reproduce the work (e.g. give exact lambda spacing, form and parameters of soft core potentials, etc.). 
One specific thing I noticed is that it is unclear if the free energy calculations involve independent lambda values or are HREX simulations. 
This has major implications for understanding their results, so it is critical to explain.\label{comment_11}
\end{point}

% Our reply
\begin{reply}
	We thank the reviewer for pointing this out. 
   We have added more details in the revised manuscript, 
   including the $\lambda$-spacing 
   (24 $\lambda$-windows for charge-change and core-hopping transformations, 
   16 $\lambda$-windows for other transformation types for RBFE; 
   32 $\lambda$-windows for neutral ligands and 48 $\lambda$-windows for charged ligands for ABFE), 
   the soft-core potential parameters, and we have clarified that all free energy calculations employ independent $\lambda$-windows rather than Hamiltonian replica exchange (HREX) simulations.
\end{reply}

\begin{revision}
   Add more details in manuscript, and paste the changes here for the reviewers' convenience.
\end{revision}

\begin{point}
Clearly, water swap MC can achieve at least two things: (a) better equilibration of waters, and (b) better sampling of water states. 
Here, in most experiments, both are mixed together -- particularly, results are looking at improved water sampling IN free energy calculations. 
How much of the present gains could have been achieved simply by better equilibration? 
e.g. with simulations at independent lambda values, 
one could use water swap MC only during the equilibration phase to better equilibrate the water, 
then turn it off during the free energy calculations. 
How much would that have improved results and what would be the relative cost/benefit? 
(Even Schrodinger's work on GCMC relative free energy calculations, 
if read very carefully, highlights that most of the benefit from their GCMC algorithm actually comes from better equilbiration.)\label{comment_12}
\end{point}

% Ye: We should do this experiment and report the results.
% Our reply
\begin{reply}
We thank the reviewer for this insightful comment, which touches on an important mechanistic question. 
To disentangle the contributions of equilibration and production-phase sampling, 
we have performed additional control simulations in which SwapMC was applied exclusively during the equilibration phase and then disabled for the production FEP windows. 
The results, presented in the revised manuscript, show that equilibration-only SwapMC recovers a portion of the accuracy gains observed with full SwapMC, 
consistent with the reviewer's observation regarding Schrödinger's GCMC work. 
We have added a discussion of these findings and the corresponding cost/benefit analysis in the revised manuscript.
\end{reply}

\begin{revision}
Add results from control simulations with SwapMC applied only during equilibration, and add discussion of the findings in the revised manuscript. \\
And paste the changes here for the reviewers' convenience.
\end{revision}


\begin{point}
What about in systems where a binding site region partially collapses when water is absent? 
e.g. in Reference 34 a key finding is that in some cases, binding sites collapse and are hard to refill with water once a water is removed. 
Presumably, this approach would have a hard time dealing with such cases 
(as did the work in reference 34, but it might be even worse here); this should be addressed.\label{comment_13}
\end{point}

% Our reply
\begin{reply}
	TBD
\end{reply}

\begin{point}
It should be disclosed whether/where this code is available, if at all. 
It may be that this is only available in a commercial package (is Uni-FEP commercial?) and if so this should be stated. 
(And if that's the case, then the authors should give much more detail to allow reproduction of the work, 
see point 1 above, otherwise this would essentially just be a software advertisement.)\label{comment_14}
\end{point}

% Our reply
\begin{reply}
	Thanks for the comment.
   We have added a statement in the revised manuscript that the code for the water swap MC method is implemented in the Uni-FEP software, which is a commercial package. 
   We have also provided more details in the methods section to allow for reproduction of the work, as mentioned in our response to comment 1.
\end{reply}

\begin{point}
Probably also cite OpenFF in connection with recent developments on force fields; in comparison with CGenFF and GAFF2, OpenFF has been more active in recent years. \label{comment_15}
\end{point}

% Our reply
\begin{reply}
	%Thanks for the suggestion. 
   %We have added the citation for OpenFF papers~\cite{Boothroyd2023} in the revised manuscript when discussing recent developments on force fields.
   We thank the reviewer for this suggestion. We have added the citation for OpenFF~\cite{Boothroyd2023, wang2024open} in the Introduction where recent force field developments are discussed, alongside CGenFF and GAFF2.
\end{reply}

\begin{point}
It would be helpful to try and make the paper at least reasonably self-contained for those from outside the area, 
e.g. by adding single-sentence explanations when relevant. 
Concrete example: p11 line 11, "the Gumbel-Max Trick" is explained in terms of "Gumbel noise" which is unhelpful to me since I know neither the trick nor what Gumbel noise is. \label{comment_16}
\end{point}

% Our reply
\begin{reply}
We thank the reviewer for this comment. 
We have revised the relevant passage in the Methods section (Implementation of Swap Monte Carlo in OpenMM) to provide a self-contained explanation. The revised text now reads:
\end{reply}

\begin{revision}
   Rememeber to paste the revised text here for the reviewers' convenience.
\end{revision}

\begin{point}
Concerning efficient evaluations of move proposals, 
could further optimizations be possible by pre-filtering candidate sites for obvious steric clashes (which would eliminate the need to do a full energy evaluation)? 
Or are the energy evaluations so efficient that this wouldn't provide further gains? 
(It would be good to briefly mention this in the paper.) \label{comment_17}
\end{point}

% Our reply
\begin{reply}

Rewrite this to indicate the randomness related with detailed balance.

We thank the reviewer for this thoughtful suggestion. 
In the current GPU-parallelized implementation, all 18,000 candidate insertion conformations are evaluated simultaneously in a single GPU kernel, 
such that the wall-clock time of the energy evaluation is largely independent of the number of candidates screened. 
As a result, pre-filtering based on steric clash detection would not reduce the per-attempt computational cost in the current architecture, although it could in principle reduce memory bandwidth requirements. We have added a brief discussion of this point in the Implementation section, noting that steric pre-filtering remains a potentially useful optimization for CPU-based or memory-bandwidth-limited implementations, and we flag it as a direction for future engineering work.
\end{reply}

\begin{point}
page 18 line 24 and vicinity: Lack of improvement in some cases is blamed on force field accuracy, protein conformational sampling. Probably other example issues should also be listed and it should be made more clear that this is just a set of examples; for example, the benchmarking work of Hahn et al. has highlighted that in many cases, ligand pose quality and system prep (selection of protonation states etc.) are a dominant source of error, and these factors aren't even mentioned here. \label{comment_18}
\end{point}

% Our reply
\begin{reply}
	TBD
\end{reply}


\newpage
%%%%%%%%%%%%%%%%%%%%%%%%%%%%%%%%%%%%%%%%%%%%%%%%%%%
\reviewersection
The authors reported an MC method for sampling water molecules in buried pockets 
that appears to help the convergence of free energy calculations for such systems. 
The topic is of interest to the readership of JCTC and results are interesting. 
The authors need to address the following few comments before the paper can be published:

%%%%%%%%%%%%%%%%
\begin{point}
While the formulas appear correct for the MC water swap, 
the parallel evaluations of multiple MC moves might break the detailed balance condition. 
In particular, after all the potential insertion conformations are sequentially evaluated, 
"the first conformation to satisfy the Metropolis criterion (random number $\le$ ASi→Sj ) is accepted, and the iteration proceeds to Step 5." 
Is there a particular fixed sequence about the different positions and orientations of the candidate water configuration for insertion? 
Or is it random as well? 
What does 'first' here mean? "first" in calculation order that satisfies the MC acceptance or 'first' in the predetermined order? 
Depending on the implementation details, the detailed balance condition may break. \label{comment_21}
\end{point}

\begin{reply}
We thank the reviewer for this careful and perceptive comment. The reviewer correctly identifies a subtlety that requires clarification.

In our implementation, the 18,000 candidate insertion conformations per swap attempt 
are generated by combining 300 uniformly sampled insertion sites with 60 approximately uniformly distributed orientations (sampled via quaternion parameterization of SO(3)). 
These conformations are re-generated independently for each swap attempt, 
so the sequence is effectively random across attempts. 
Within a single attempt, conformations are evaluated in a fixed order, 
and "first" in Step 4 refers to the first conformation in this order that satisfies the Metropolis acceptance criterion.
Since the conformation set is independently resampled at each attempt, the insertion probability $g^{ins}_{S_j}$ remains uniform over the destination region, and detailed balance is preserved as derived in Eq. 7.

We have substantially expanded the relevant section in the Methods to make this point explicit and to clarify the implementation details, 
including the fixed ordering and the independence of each acceptance test.
\end{reply}

\begin{revision}
Make revision on the manuscript to clarify the implementation details of the parallel MC move evaluations and how detailed balance is preserved, and paste the changes here for the reviewers' convenience.
\end{revision}


%%%%%%%%%%%%%%%%
\begin{point}
The validation on free energy calculation is good, but not quite sufficient. 
I recommend control simulations on more basic systems to validate the theory and implementation. 
For example, the authors can perform simulations on ideal gas systems where analytical solutions are available, 
and a box of water molecules where water exchange is not a problem to compare the simulation results with NPT without MC swap.\label{comment_22}
\end{point}

\begin{reply}
We thank the reviewer for this suggestion. 
We respectfully note that the SwapMC method is specifically designed for water exchange between a defined cavity subspace a
nd the surrounding bulk within a periodic boundary condition (PBC) box, 
and its implementation is tightly integrated with the alchemical FEP workflow in Uni-FEP. 
Constructing standalone validation systems (ideal gas or pure bulk water boxes) 
would require substantial re-engineering outside this framework and is therefore not straightforward to execute.
\newline
\noindent More fundamentally, we argue that the theoretical correctness of the acceptance criterion is established analytically through the detailed balance derivation in Eqs. 2–8, 
which holds for any partition of a PBC simulation box into two complementary regions. 
The practical validity of the implementation is then demonstrated by the robustness of the results across 94 RBFE and 67 ABFE calculations: SwapMC consistently drives cavity water occupancy toward physically reasonable values (Fig. 2), 
improves thermodynamic accuracy across diverse protein targets (Tables 1–2), 
and yields results insensitive to initial water configuration — behavior that would not be expected from an incorrectly implemented sampler. 
We therefore consider the existing benchmarks to constitute sufficient empirical validation of both the theory and the implementation. 
We also note that we plan to refactor SwapMC into a standalone, framework-independent water sampling module in future work, 
which will facilitate more straightforward validation against simple reference systems and broader adoption by the community.
\end{reply}

%%%%%%%%%%%%%%%%
\begin{point}
The analysis of AbFEP results needs improvement. 
Experimental dG has contributions from interaction free energy with the ligand and the apo to holo protein conformational change, 
with the later missing in short AbFEP simulations. 
Instead of comparing the raw predicted dG with exp, the authors are encouraged to separate out the systematic shift as compared to exp (related to apo to holo conformational free energy missing for all ligands), 
and then access the effort of MC swap on both the systematic error and the ranking order error. \label{comment_23}
\end{point}

\begin{reply}
We thank the reviewer for this valuable suggestion. 
The reviewer correctly points out that the systematic offset between predicted and experimental absolute binding free energies 
— arising from incomplete sampling of the apo-to-holo conformational change — affects all ligands within a given target approximately equally. 
We have revised the ABFE analysis accordingly: for each protein target, we subtract the mean signed error (MSE) across all ligands to remove this systematic contribution and report the resulting shifted RMSE. 
We additionally report Kendall's ranking value $\tau$ to evaluate ranking accuracy independently of absolute errors. 
Table 2 and the surrounding discussion have been updated to reflect this analysis.
\end{reply}

\begin{revision}
   Paste the revision of the ABFE analysis and discussion here for the reviewers' convenience.
\end{revision}

%%%%%%%%%%%%%%%%
\begin{point}
The raw data for both RbFEP and AbFEP need to be included in the SI.\label{comment_24}
\end{point}

\begin{reply}
We thank the reviewer for this request. 
We have included the complete raw data for all RBFE and ABFE calculations as Supporting Information. 
Specifically, the SI now contains: (1) per-edge $\Delta \Delta G$ predictions and experimental values for all 94 RBFE transformations across the eight water-displacement benchmark systems; 
and (2) per-ligand $\Delta G$ predictions and experimental values for all 67 ABFE calculations across the eight JACS benchmark targets. 
Both with-SwapMC and without-SwapMC results are provided for all entries.
\end{reply}

\bibliography{ref,common_ref}
\bibliographystyle{chicago}
\end{document}
